% !PRINTED 202
\weberpage{202}

Hieraus folgt nun, daß man die Teilung durch 2 und mithin
auch durch jede Potenz von 2 durch eine Kette von Quadrat\-
wurzeln ausführen kann. Setzt man daher die Aufgabe der Tei\-
lung durch eine ungerade Zahl als gelöst voraus, so ist die
Teilung durch eine gerade Zahl auf Quadratwurzeln zurückgeführt. Im
folgenden beschäftigen wir uns ausschließlich mit \gesperrt{der Teilung
durch ungerade Zahlen}.

\begin{center}
§ 60. \gesperrt{Die Teilung durch eine ungerade Zahl.}
\end{center}

Setzt man, wenn $n$ eine ungerade Zahl ist,

\weq{1}{x = \mathrm{sn}\,\frac{v}{n}, \qquad y = \mathrm{cn}\,\frac{v}{n}, \qquad z = \mathrm{dn}\,\frac{v}{n},}

so erhält man aus den Gleichungen IV, \S\ 57:

\weq{2}{\begin{aligned}
D(x^2)\,\mathrm{sn}\,v - x\,A(x^2) &= 0,\\
D(x^2)\,\mathrm{cn}\,v - y\,B(x^2) &= 0,\\
D(x^2)\,\mathrm{dn}\,v - z\,C(x^2) &= 0.
\end{aligned}}

Die erste dieser Gleichungen ist in bezug auf die Unbekannte
$x$ vom Grade $n^2$; durch die zweite und dritte werden $y$ und $z$
\gesperrt{rational} durch $x$ (und durch cn\,$v$, dn\,$v$) ausgedrückt. Es ergeben
sich also $n^2$ Wertsysteme für die drei Unbekannten $x, y, z$, deren
Bedeutung ist:

\weq{3}{\begin{aligned}
x_{\mu,\mu'} &= \mathrm{sn}\left(\frac{v}{n} + \frac{4\mu K+4\mu' iK'}{n}\right),\\
y_{\mu,\mu'} &= \mathrm{cn}\left(\frac{v}{n} + \frac{4\mu K+4\mu' iK'}{n}\right),\\
z_{\mu,\mu'} &= \mathrm{dn}\left(\frac{v}{n} + \frac{4\mu K+4\mu' iK'}{n}\right),
\end{aligned}}

worin $\mu, \mu'$ je ein vollständiges Restsystem $(\mathrm{mod}\ n)$
durchlaufen.

Die erste Gleichung (2):

\weq{4}{D(x^2)\,\mathrm{sn}\,v - x\,A(x^2) = 0,}

vom Grade $n^2$ heißt die \gesperrt{allgemeine Teilungsgleichung}. Sie
ist in dem Sinne irreducibel, daß sie nicht in Faktoren zerlegbar
ist, die in bezug auf sn\,$v$, cn\,$v$, dn\,$v$ rational sind und beliebige
von $v$ unabhängige Koeffizienten haben.

Denn zunächst sind die $n^2$ Werte $x_{\mu,\mu'}$ alle voneinander ver\-
schieden, und wenn irgend eine rationale Gleichung:

\[
F\left(\mathrm{sn}\,\frac{v}{n},\ \mathrm{sn}\,v,\ \mathrm{cn}\,v,\ \mathrm{dn}\,v\right) = 0
\]
%CONT>>
